\makeatletter
\def\input@path{{./}{../}{../../}{../../../}{../../../../}{preamble/}{../preamble/}{../../preamble/}{../../../preamble/}}
\makeatother
\input{preamble_exams.tex}
\input{graphs/highlights.tex}
% \GTenHighlightsfalse % Uncomment to disable instructional highlights.
\input{graphs/specific_highlights.tex}

\hypersetup{pdftitle={Grade 10 Global Economics: Term 3, Lesson 1, C9},pdfsubject={Home practice: Two-Asset Portfolios and Diversification}}
\title{Grade 10 Global Economics\\Term 3, Lesson 1: Home Practice\\C9 --- Two-Asset Portfolios and Diversification}
\author{}
\date{}
\begin{document}
\maketitle
\setcounter{tocdepth}{1}
\setcounter{secnumdepth}{2}
\phantomsection\label{toc}
\tableofcontents
\section*{Evaluated criteria}
\begin{description}
\item[\textbf{C9}] Analyses investment portfolios by interpreting assets, weights, expected returns, variance, covariance, and correlation as elements that determine portfolio risk.
\end{description}
\textbf{How to practise.} Keep annual returns in percentage-point units throughout. Write the weights and their sum before using the portfolio formulas, and retain the covariance term in every variance calculation. Try each task on paper before opening its solution link. Record the model, working, units and an interpretation; use the arrows to return.
All datasets and economic or financial models in this activity are illustrative, rather than observed market evidence.

\clearpage
\section{Introduction}\label{sec:introduction}
\subsection*{Highlight legend}
\phantomsection\label{tab:highlight-legend}
\begin{tabularx}{\linewidth}{@{}>{\raggedright\arraybackslash}p{.20\linewidth}>{\raggedright\arraybackslash}p{.22\linewidth}>{\raggedright\arraybackslash}X@{}}
\toprule
\textbf{Symbol or notation} & \textbf{Concept} & \textbf{Brief meaning}\\
\midrule
$\GTenHLParameter{w_A,w_B}$ & \GTenHLParameter{Weights} & Fractions of starting capital. \\
$R_p=w_AR_A+w_BR_B$ & Portfolio return & Weighted asset return. \\
$\HLExpectedValue{\mu_A,\mu_B,\mu_p}$ & \HLExpectedValue{Expected returns} & Model average returns. \\
$\HLPopulationSD{\sigma_i},\ \HLPopulationVariance{\sigma_i^2}$ & \HLPopulationSD{SD} and \HLPopulationVariance{variance} & Individual return spread. \\
$\GTenHLSecond{\operatorname{Cov}(R_A,R_B)}$ & \GTenHLSecond{Covariance} & Joint movement in squared units. \\
$\GTenHLThird{\rho}$ & \GTenHLThird{Correlation} & Unitless direction and strength. \\
\bottomrule
\end{tabularx}
\subsection{Introductory example 1: Dividing an investment between two assets}\label{ex:intro-one}
A portfolio allocates 60 dollars of each 100 dollars to asset A and 40 to asset B. Its \GTenHLParameter{weights} are $\GTenHLParameter{w_A=0.60}$ and $\GTenHLParameter{w_B=0.40}$, summing to one. If annual expected returns are $\HLExpectedValue{6}$ and $\HLExpectedValue{10}$ percentage points, the portfolio expected return is $0.60(6)+0.40(10)=\HLExpectedValue{7.6}$ percentage points, meaning an expected $7.6\%$ gain.
\subsection{Introductory example 2: Assets can move together}\label{ex:intro-two}
Suppose A and B have annual SDs $\HLPopulationSD{10}$ and $\HLPopulationSD{20}$ percentage points and \GTenHLThird{correlation} $\GTenHLThird{\rho=0.25}$. Their \GTenHLSecond{covariance} is $\GTenHLSecond{0.25(10)(20)=50}$ squared percentage points. Positive covariance adds to portfolio variance; changing correlation changes risk even if both assets' separate means and SDs remain fixed.
\section{Notation Reference and Practice Procedure}\label{sec:notation-and-procedure}
\subsection{Notation reference}\label{tab:notation-reference}
\begin{tabularx}{\linewidth}{@{}>{\raggedright\arraybackslash}p{.16\linewidth}>{\raggedright\arraybackslash}p{.18\linewidth}>{\raggedright\arraybackslash}p{.27\linewidth}>{\raggedright\arraybackslash}X@{}}
\toprule
\textbf{Symbol or notation} & \textbf{Concept} & \textbf{Definition} & \textbf{Meaning in the introductory examples}\\
\midrule
$\GTenHLParameter{w_A,w_B}$ & \GTenHLParameter{Weights} & Here $0le w_A,w_Ble1$ and $w_A+w_B=1$; no borrowing or short selling. & \textbullet\ \hyperref[ex:intro-one]{Example 1}: The allocation fractions are 0.60 and 0.40. \\
$R_p=w_AR_A+w_BR_B$ & Portfolio return & For a common horizon, portfolio return is the capital-weighted sum of asset returns. & \textbullet\ \hyperref[ex:intro-one]{Example 1}: The two annual assets make one annual portfolio. \\
$\HLExpectedValue{\mu_A,\mu_B,\mu_p}$ & \HLExpectedValue{Expected returns} & $mu_p=w_Amu_A+w_Bmu_B$; expectation does not require independence. & \textbullet\ \hyperref[ex:intro-one]{Example 1}: The expected return is 7.6 percentage points. \\
$\HLPopulationSD{\sigma_i},\ \HLPopulationVariance{\sigma_i^2}$ & \HLPopulationSD{SD} and \HLPopulationVariance{variance} & SD has return units; variance has squared return units. & \textbullet\ \hyperref[ex:intro-two]{Example 2}: SDs are 10 and 20 percentage points. \\
$\GTenHLSecond{\operatorname{Cov}(R_A,R_B)}$ & \GTenHLSecond{Covariance} & $E[(R_A-mu_A)(R_B-mu_B)]$; equals $
hosigma_Asigma_B$. & \textbullet\ \hyperref[ex:intro-two]{Example 2}: Covariance is 50 squared percentage points. \\
$\GTenHLThird{\rho}$ & \GTenHLThird{Correlation} & Covariance divided by $sigma_Asigma_B$ for positive SDs; lies in $[-1,1]$. & \textbullet\ \hyperref[ex:intro-two]{Example 2}: Correlation 0.25 indicates positive linear co-movement. \\
\bottomrule
\end{tabularx}
\subsection{Practice procedure}
\begin{enumerate}
\item List both assets, a common horizon, expected returns and SDs. State whether numeric returns are decimals or percentage points.
\item Convert capital amounts to weights and confirm they sum to one. Compute $\mu_p=w_A\mu_A+w_B\mu_B$.
\item Use $\operatorname{Cov}(R_A,R_B)=\rho\sigma_A\sigma_B$ when needed. Calculate both weighted variance terms and the covariance term.
\item Take the square root of portfolio variance for SD. Compare expected return and risk together under any stated constraints; diversification can reduce risk but does not generally remove it.
\end{enumerate}
\subsection{Portfolio formulas and return units}
For two assets held over the same year, without rebalancing within the year,
\[
\begin{aligned}
R_p&=w_AR_A+w_BR_B,\\
\mu_p&=w_A\mu_A+w_B\mu_B,\\
\sigma_p^2&=w_A^2\sigma_A^2+w_B^2\sigma_B^2
       +2w_Aw_B\operatorname{Cov}(R_A,R_B),\\
\operatorname{Cov}(R_A,R_B)&=\rho\sigma_A\sigma_B,\qquad
\sigma_p=\sqrt{\sigma_p^2}.
\end{aligned}
\]
Here numeric return $6$ means $6\%$, so means and SDs use \textbf{percentage points} and variances and covariances use \textbf{squared percentage points}. For initial value $V$ dollars, monetary gain is $VR_p/100$ dollars. Decimal-return inputs would require dividing means and SDs by 100 and variances and covariances by $10{,}000$. Do not mix the two systems. Covariance is zero when assets are independent, but zero covariance alone does not establish independence.
\section{General Worked Problems}\label{sec:general-worked}
\SectionProblemsTableOfContents{
\SectionProblemLink{prob:sec:general-worked0}{Capital weights and expected gain}\\
\SectionProblemLink{prob:sec:general-worked1}{Portfolio variance with zero covariance}\\
\SectionProblemLink{prob:sec:general-worked2}{Adding positive covariance}}

\subsection{Capital weights and expected gain}\label{prob:sec:general-worked0}
An investor allocates 600 dollars to asset A and 400 dollars to asset B for one year. Their expected annual returns are 6 and 10 percentage points. No contributions, withdrawals or borrowing occur.
\QuestionsTasks[sec:general-worked]{{9}/{sol:sec:general-worked0}/{Calculate the portfolio weights, expected annual return and expected dollar gain. Compare with an equal-dollar allocation and explain what cannot yet be concluded about risk.}}
\subsubsection*{Worked solution}\label{sol:sec:general-worked0}
The total starting value is $1000$ dollars. Weights are $w_A=600/1000=\GTenHLParameter{0.60}$ and $w_B=400/1000=\GTenHLParameter{0.40}$, with sum one.
\[
\mu_p=0.60(6)+0.40(10)=\HLExpectedValue{7.6}
\]
percentage points, so expected gain is $1000(7.6/100)=76$ dollars. With equal-dollar weights,
\[
\mu_p=0.50(6)+0.50(10)=\HLExpectedValue{8},
\]
giving expected gain $80$ dollars on the same capital. The expected gain rises by four dollars. These are averages under the illustrative model, not promised gains. We cannot compare risk from expected returns alone: asset variances and their covariance are also needed.
\secInicio[prob:sec:general-worked0][sec:general-worked]


\subsection{Portfolio variance with zero covariance}\label{prob:sec:general-worked1}
Two assets have expected annual returns $\mu_A=6$, $\mu_B=10$ and annual SDs $\sigma_A=10$, $\sigma_B=20$, all in percentage points. Their covariance is zero. Consider equal weights $w_A=w_B=0.50$.
\QuestionsTasks[sec:general-worked]{{9}/{sol:sec:general-worked1}/{Calculate the portfolio mean, variance and SD. Compare its SD with the weighted average of the asset SDs, explaining the distinction.}}
\subsubsection*{Worked solution}\label{sol:sec:general-worked1}
The expected return is $\mu_p=0.5(6)+0.5(10)=\HLExpectedValue{8}$ percentage points. The risk calculation uses squared weights:
\[
\begin{aligned}
\sigma_p^2
&=(0.5)^2(10^2)+(0.5)^2(20^2)+2(0.5)(0.5)(0)\\
&=25+100+0=\HLPopulationVariance{125}
\end{aligned}
\]
squared percentage points. Thus $\sigma_p=\sqrt{125}\approx\HLPopulationSD{11.18}$ percentage points.
The weighted average of the separate SDs is $0.5(10)+0.5(20)=15$ percentage points, which is not the portfolio SD here. Combining imperfectly co-moving assets reduces portfolio risk relative to that benchmark. It does not make the portfolio safer than every component: $11.18$ is still above A's SD of 10, and risk remains positive.
\secInicio[prob:sec:general-worked1][sec:general-worked]


\subsection{Adding positive covariance}\label{prob:sec:general-worked2}
Keep the preceding annual asset inputs: means 6 and 10, SDs 10 and 20, and equal weights. Change their correlation from zero to $\rho=0.25$.
\QuestionsTasks[sec:general-worked]{{9}/{sol:sec:general-worked2}/{Find covariance and recalculate portfolio risk. Explain which results change and which remain fixed.}}
\subsubsection*{Worked solution}\label{sol:sec:general-worked2}
Correlation is unitless. Convert it to covariance:
\[
\operatorname{Cov}(R_A,R_B)=0.25(10)(20)=\GTenHLSecond{50}
\]
squared percentage points. The portfolio variance is
\[
\sigma_p^2=25+100+2(0.5)(0.5)(50)
=125+25=\HLPopulationVariance{150}.
\]
Its SD is $\sqrt{150}\approx\HLPopulationSD{12.25}$ percentage points, exceeding the zero-covariance value of about $11.18$. Positive co-movement increases the joint-risk contribution. The mean stays $0.5(6)+0.5(10)=8$ percentage points because weights and asset means are unchanged. The portfolio still has lower SD than the weighted-average benchmark 15, but does not eliminate risk.
\secInicio[prob:sec:general-worked2][sec:general-worked]
\section{Advanced Problems}\label{sec:advanced}
\SectionProblemsTableOfContents{
\SectionProblemLink{prob:sec:advanced0}{Correlation sensitivity with fixed assets}\\
\SectionProblemLink{prob:sec:advanced1}{Choosing under a return target and risk limit}}

\subsection{Correlation sensitivity with fixed assets}\label{prob:sec:advanced0}
Assets A and B have annual means 6 and 10 percentage points and SDs 10 and 20 percentage points. Fix equal weights and compare correlations $-0.5$, $0$, $0.5$ and $1$.
\QuestionsTasks[sec:advanced]{{9}/{sol:sec:advanced0}/{Calculate the variance and SD under each assumption. Rank the risks and explain diversification without assuming that holding two assets eliminates uncertainty.}}
\subsubsection*{Worked solution}\label{sol:sec:advanced0}
The weighted individual variance terms total $25+100=125$. Since covariance is $200\rho$, the covariance contribution is $2(0.5)(0.5)(200\rho)=100\rho$. Therefore
\[
\sigma_p^2=125+100\rho.
\]
\[
\begin{array}{c|rrrr}
\rho&-0.5&0&0.5&1\\ \hline
\operatorname{Cov}\text{ (squared pp)}&-100&0&100&200\\
\sigma_p^2\text{ (squared pp)}&75&125&175&225\\
\sigma_p\text{ (pp), approximately}&8.66&11.18&13.23&15.00
\end{array}
\]
Risk increases with correlation for these positive fixed weights. At correlation one, SD equals the weighted average $15$: the assets move together perfectly and this combination has no reduction relative to that benchmark. At $-0.5$, offsetting movements reduce SD to about $8.66$, but variance remains positive. All four cases retain mean eight percentage points. Diversification affects dispersion through joint movement; common shocks, uncertain future correlations and remaining asset variability can leave substantial risk.
\secInicio[prob:sec:advanced0][sec:advanced]


\subsection{Choosing under a return target and risk limit}\label{prob:sec:advanced1}
Use annual means $\mu_A=6,\mu_B=10$ and SDs $\sigma_A=10,\sigma_B=20$ percentage points. Candidate P has weights $(0.50,0.50)$ and candidate Q has weights $(0.80,0.20)$. First assume zero covariance. A policy requires expected return at least 8 percentage points and SD at most 12 percentage points. Then revise correlation to $0.25$.
\QuestionsTasks[sec:advanced]{{9}/{sol:sec:advanced1}/{Compare both candidates under the policy, then explain whether the choice remains feasible after the correlation revision.}}
\subsubsection*{Worked solution}\label{sol:sec:advanced1}
Under zero covariance, P has mean $8$ and variance $125$, hence SD about $11.18$. It satisfies both limits. For Q,
\[
\mu_Q=0.8(6)+0.2(10)=\HLExpectedValue{6.8},
\]
\[
\sigma_Q^2=0.8^2(10^2)+0.2^2(20^2)=64+16=\HLPopulationVariance{80},
\quad \sigma_Q\approx8.94.
\]
Q has lower SD but misses the expected-return target, so P is the only feasible one of these two candidates.
With correlation $0.25$, covariance becomes $50$ squared percentage points. P's variance rises to $150$ and SD to about $12.25$, above the risk cap. Q's covariance term is $2(0.8)(0.2)(50)=16$, so its variance becomes $96$ and SD about $9.80$. Its mean remains $6.8$, still below target. Thus neither listed candidate is feasible under the revised joint-risk assumption. This does not prove that every possible weight choice fails; it establishes the result only for P and Q.
\secInicio[prob:sec:advanced1][sec:advanced]

\end{document}

